\chapter{Introduction}


%% Definitely need to explain this section better, but it serves as motivation for what we are doing
A key area of research in modern algebraic number theory are the so called eigenvarieties, these
are rigid analytic spaces that lie above a weight space, whose points are isomorphic to
generalised eigenvectors of Hecke operators.

A point in weight space is a character, and its fiber will be family of forms of weight given by the
character. The isomorphism comes from the points in the eigenvariety being characters, which
determine the eigenvalues for the form with respect to the $T_n$ operator. Fixing a specific operator,
say $U_p$, we can look at the fiber of a character and then get a family of eigenvectors for $U_p$.

Formalising the required eigenvariety machinery is out of scope for this thesis. However, studying
what the family of eigenforms looks like for a given character is an interesting question. It turns
out this question can have surprising answers - and our goal of this thesis is to formalise
such an answer.

To do this we will restrict ourselves to quaternionic modular forms over definite quaternion algebras.
Jacquet-Langlands suggests that these should relate to certain families of classical modular forms.
However, compared to the analytical nature of the classical forms, quaternionic modular forms are
much nicer to work with, being combinatorial in nature.

The aim is to define quaternionic modular forms and their Hecke operators in Lean, and then following
\cite{Jacobs} show that the slopes (eigenvalues) of compact Hecke operators are in arithmetic
progression. That is, without defining eigenvarieties, we are able to discuss what the geometry
if a fiber looks like in the "quaternionic" eigenvariety.


To prove these statements we need to formalise a lot of mathematics from the ground up. This blueprint
will consist of numerous distinct sections that we need to ultimately formalise our goal. A rough
outline of what this looks like is given below:

%% I want to generate a flow chart of what the dependencies look like... for now just writing them as a list
%% This could be generated by my blueprint eventually... but think I want a basic version first

\begin{enumerate}
  \item Restricted power series and the gauss norm
  \item Weierstrass preparation
  \item Newton polygons
  \item Nonarchimedean compact operators and Fredholm theory
  \item Quaternionic modular forms and Fujisakis lemma
  \item Hecke operators
  \item the main result
\end{enumerate}
